\chapter{Motor digital}
\label{ch:ch4_engine}

% [BORRADOR v1] Material de partida: docs/arquitectura.md (sec. 9), docs/mapa_registros.md,
% docs/estado_del_arte_cripto.md, rtl/ (aes, drbg, trng, io, top).

El motor digital convierte los bits en bruto de la fuente de entropía en
claves y las usa, y está escrito de forma que el mismo \ac{RTL} funcione
en la FPGA usada para la verificación y, más adelante, en celdas
estándar. Nada en él depende de un fabricante: ni primitivas, ni macros
de memoria, ni bloques de gestión de reloj. El único bloque dependiente
de la tecnología en todo el diseño es la pareja de osciladores en anillo
del Capítulo~\ref{ch:ch3_entropy}.

\section{Arquitectura y la decisión de una única primitiva}
\label{sec:arch}

La Figura~\ref{fig:engine} resume la ruta de datos. Los bits en bruto
salen del \ac{ERO} en el dominio de reloj del anillo de muestreo, cruzan
al reloj del sistema de bit en bit y alimentan a tres consumidores: los
tests de salud, el colector de semilla del generador determinista y un
búfer de captura que solo se usa para la caracterización. Las claves
salen de un \ac{CTRDRBG}; la clave actual se carga en un bloque AES-CMAC
que responde a retos; y un esclavo \ac{I2C} lo expone todo como un
mapa de registros.

\begin{figure}[htbp]
\centering
\resizebox{\textwidth}{!}{%
\begin{tikzpicture}[node distance=6mm and 8mm, >=Stealth,
  bloque/.style={draw, rounded corners=2pt, minimum height=8mm, minimum width=19mm, align=center, font=\small},
  flecha/.style={->, thick}]
\node[bloque] (ro) {RO1, RO2\\(LUT / transistor)};
\node[bloque, right=of ro] (ero) {ERO\\muestreador $\div K_D$};
\node[bloque, right=of ero] (cdc) {CDC de bit};
\node[bloque, above right=of cdc] (ht) {Tests de salud\\RCT, APT};
\node[bloque, right=of cdc] (seed) {Colector de semilla\\384 bits};
\node[bloque, below right=of cdc] (cap) {Captura\\4096 bits};
\node[bloque, right=of seed] (drbg) {CTR\_DRBG\\+ df};
\node[bloque, right=of drbg] (cmac) {AES-CMAC};
\node[bloque, below=12mm of drbg] (i2c) {Esclavo I2C\\mapa de registros};
\draw[flecha] (ro) -- (ero); \draw[flecha] (ero) -- (cdc);
\draw[flecha] (cdc) -- (ht); \draw[flecha] (cdc) -- (seed);
\draw[flecha] (ero) |- (cap);
\draw[flecha] (seed) -- (drbg) node[midway, above, font=\scriptsize] {entropía, nonce};
\draw[flecha] (drbg) -- (cmac) node[midway, above, font=\scriptsize] {clave};
\draw[flecha, <->] (i2c) -- (drbg); \draw[flecha, <->] (i2c) -| (cmac);
\draw[flecha, <->] (i2c) -| (cap); \draw[flecha] (ht) -| (i2c);
\end{tikzpicture}}
\caption{Ruta de datos del motor. Solo los osciladores en anillo dependen
de la tecnología; todos los demás bloques son RTL portable.}
\label{fig:engine}
\end{figure}

La decisión que da forma a todo el motor es usar \textbf{una sola
primitiva criptográfica, AES-128, para las tres tareas}: acondicionar la
entropía, generar claves y autenticar. La primera arquitectura tenía un
núcleo SHA-256 como acondicionador; las estimaciones de síntesis del
Capítulo~\ref{ch:ch6_asic} mostraron que ocuparía casi la mitad del área
digital, y SP~800-90B incluye construcciones basadas en AES --- CMAC,
CBC-MAC y la función de derivación basada en cifrado de bloque de
SP~800-90A --- entre sus componentes de acondicionamiento
\emph{aprobados}~\cite{turan_2018_sp80090b}. Prescindir del hash lleva el
motor de un valor estimado de $\sim$25 a $\sim$13~\acp{GE} (miles), a
costa de pasar de un \ac{CTRDRBG} \emph{sin} función de derivación, que
necesita una entrada con entropía completa, a uno \emph{con} ella, que es
la configuración habitual en la práctica y admite directamente entradas
sin entropía completa. El núcleo AES solo cifra: el modo contador y CMAC
nunca descifran, y ofrecer el descifrado como servicio al ESP32 no
tendría sentido, ya que el ESP32 tiene su propio acelerador. La
verificación cruzada no se ve afectada: el motor cifra y el ESP32
descifra con mbedtls.

\section{Núcleo de cifrado AES-128}
\label{sec:aes}

El núcleo es iterativo, una ronda por ciclo de reloj, once ciclos por
bloque (suma inicial de clave más diez rondas), con las claves de ronda
expandidas al vuelo a partir de la anterior --- el cifrado solo necesita
la expansión hacia delante, y así se evita almacenar once claves de
ronda. A 100~MHz eso supone 1,16~Gbit/s de tasa de bits del núcleo, muy
por encima de lo que puede transportar cualquier enlace \ac{I2C}, así que
segmentar el núcleo (pipelining) sería área desperdiciada; en su lugar se
elige la forma iterativa porque una ronda es una función combinacional
pura que puede comprobarse directamente contra la
norma~\cite{nist_2023_fips197}. La S-box es una tabla de 256 entradas
sintetizada como lógica, MixColumns se escribe como la red \texttt{xtime}
habitual, y el estado sigue la convención por columnas de FIPS~197 con
el byte~0 a la izquierda, de modo que las cadenas hexadecimales de la
norma se corresponden una a una con los vectores VHDL.

El núcleo se verificó con los cinco vectores oficiales disponibles para
el cifrado AES-128 (apéndices~B y~C.1 de FIPS~197, y los tres bloques ECB
de SP~800-38A F.1.1) y con una cadena de mil cifrados para poner de
manifiesto cualquier estado residual entre bloques. Una comprobación que
no estaba prevista resultó ser la más convincente: el cifrado del bloque
todo ceros con la clave de ejemplo de SP~800-38A dio
\texttt{7df76b0c\ldots546f}, que es exactamente el valor intermedio $L$
publicado por el NIST en sus ejemplos de CMAC~\cite{nist_cmac_examples}.

\section{AES-CMAC}
\label{sec:cmac}

AES-CMAC~\cite{nist_2016_sp80038b, song_2006_rfc4493} es la primitiva del
reto-respuesta: el ESP32 envía un nonce y el motor devuelve su
etiqueta con la clave actual. El bloque calcula las subclaves $K_1$,
$K_2$ una sola vez al cargar una clave (un cifrado del bloque cero y dos
desplazamientos condicionales a la izquierda con la constante polinómica
\texttt{0x87}) y después ejecuta CBC-MAC sobre el mensaje, aplicando
$K_1$ si el último bloque está completo o rellenando con \texttt{0x80} y
aplicando $K_2$ en caso contrario. Se eligió frente a HMAC-SHA-256 porque
reutiliza la única primitiva y cuesta un cifrado por bloque frente a dos
compresiones del hash. El bloque supera los cuatro vectores de la
RFC~4493 (mensaje vacío, un bloque, 40 y 64 bytes), que ejercitan las dos
ramas de relleno.

\section{CTR\_DRBG con función de derivación}
\label{sec:drbg}

El generador implementa CTR\_DRBG con AES-128 y la función de derivación
basada en cifrado de bloque de SP~800-90A~\cite{barker_2015_sp80090ar1}.
Su formato de entrada es fijo a propósito: 256~bits de entropía en bruto
y 128~bits de nonce, de modo que $S = L\,\|\,N\,\|\,E\,\|\,\mathit{nonce}\,\|\,\texttt{0x80}\,\|\,0\ldots$
ocupa exactamente cuatro bloques y la función de derivación se convierte
en una secuencia cerrada de catorce cifrados; una entrada de longitud
variable costaría una ruta de datos orientada a bytes sin ningún
beneficio. La instanciación (o el resembrado) cuesta 14~operaciones AES,
183~ciclos; generar 256~bits cuesta cuatro, 53~ciclos. El estado (Key, V)
nunca se expone; el contador de resembrados es legible para que el
anfitrión pueda imponer una política de resembrado muy por debajo del
límite de $2^{48}$ de la norma.

La verificación se hizo en dos pasos. Primero, se escribió en Python un
modelo de referencia exacto al bit --- con su propio AES, independiente
del VHDL --- y se ejecutó contra los vectores del CAVP: 960 casos AES-128
con y sin función de derivación, con y sin resembrado, todos superados.
Segundo, ese modelo generó las salidas esperadas para el formato de
entrada fijo del hardware, y el VHDL las reprodujo: dos generaciones
consecutivas, un resembrado seguido de una generación, el contador de
resembrados y el hecho de que una nueva instanciación borra el estado
anterior.

\section{Tests de salud}
\label{sec:health}

Los dos tests aprobados por SP~800-90B se ejecutan de forma continua
sobre los bits en bruto~\cite{turan_2018_sp80090b}. El test de recuento
de repeticiones (\emph{Repetition Count Test}) falla cuando llegan
seguidas $C_{\mathrm{RCT}} = 1 + \lceil -\log_2\alpha / H \rceil$
muestras idénticas; el test de proporción adaptativa (\emph{Adaptive
Proportion Test}) cuenta, en una ventana de $W = 1024$ muestras (el valor
para fuentes binarias; 512 es para fuentes no binarias), cuántas repiten
la primera, y falla al llegar a
$C_{\mathrm{APT}} = 1 + \mathrm{CRITBINOM}(W, 2^{-H}, 1-\alpha)$. Para el
objetivo de diseño $H = 0{,}98$ y $\alpha = 2^{-20}$ resultan
$C_{\mathrm{RCT}} = 22$ y $C_{\mathrm{APT}} = 596$; ambos son parámetros
genéricos (generics), porque $H$ se fija a partir de la min-entropía
\emph{medida}, no se supone. El test de proporción comprueba también el
complemento de la cuenta, algo que la norma permite para fuentes
binarias: de lo contrario, una fuente atascada en el valor opuesto a la
primera muestra de la ventana puntuaría cero y pasaría. Las alarmas
quedan enclavadas hasta que se borran explícitamente, y mientras haya
alguna alarma activa no sale nada del motor: el secuenciador se niega a
cargar la semilla, a generar clave y a producir salida aleatoria. Dos
reglas más cierran el hueco entre un test continuo y una semilla de 384
bits. Primero tiene que pasar un test de arranque (SP~800-90B \S4.3):
una ventana completa del test de proporción, 1024 muestras, cerrada sin
alarma tras el reset o tras borrar las alarmas, y borrarlas reinicia la
ventana. Y una semilla solo se acepta cuando se ha cerrado sin alarma la
ventana que contiene su último bit; la semilla es más corta que la
ventana, así que de otro modo entraría en el generador antes de que el
test de proporción la hubiera juzgado.

Cuatro reglas más salieron de revisar el secuenciador contra el mapa de
registros y no contra la norma. El buffer de captura de bits en bruto
registra los mismos bits de los que se recoge la semilla, de modo que una
captura lanzada durante una siembra dejaba la semilla legible en modo de
prueba; ahora capturar es una orden de modo de prueba que se rechaza con
el motor ocupado, y la siembra se rechaza mientras hay una captura en
curso. El contador de peticiones del \ac{DRBG} ya no da la vuelta: en su
tope de 32 bits el motor se niega a generar hasta que se resiembra y lo
indica en un bit de diagnóstico, que es el límite por hardware de la
política de resembrado que el anfitrión sigue teniendo que aplicar mucho
antes. Una orden de borrado total pone a cero la clave, el estado y el
contador del generador, el reto, el registro aleatorio y la semilla
recogida, y carga la clave nula en el autenticador; se atiende incluso
con alarma, y después hay que volver a sembrar. Y el divisor de muestreo
se reinicia a su máximo y no a un valor pequeño, de modo que un
anfitrión que olvide configurarlo siembra a la tasa más lenta y con más
entropía por bit, nunca por debajo del ajuste caracterizado.

El banco de pruebas inyecta los modos de fallo para los que existen los
tests: una fuente sana supera veinte ventanas (peor racha 13, peor
proporción 540); una fuente atascada dispara el test de repeticiones
exactamente en la muestra~22; una fuente sesgada al 70\,\% dispara el
test de proporción en 713 sin disparar el test de repeticiones; y una
fuente atascada en el complemento se detecta. La primera fuente
``sana'' del banco --- un registro de desplazamiento inicializado con un
único uno --- fue detectada por los propios tests, con una racha inicial
de 27 ceros y una primera ventana al 76\,\%. El detector tenía razón y
el banco estaba mal; ahora la fuente parte de una semilla densa y se
calienta antes de usarse.

\section{Ruta de captura y cruces de dominios de reloj}
\label{sec:cdc}

Los bits en bruto se producen en el dominio de reloj de RO2 y se
consumen en el dominio del sistema. La ruta serie, bit a bit, hacia los
tests de salud y el colector de semilla usa un sincronizador por
conmutación (toggle): el origen registra el bit e invierte una bandera;
el destino sincroniza la bandera a través de tres flip-flops y lee el bit
en el flanco de la bandera. Requiere que cada bit dure al menos ocho
ciclos del sistema, es decir, $K_D \ge 64$ para un anillo de muestreo de
$\sim$700~MHz a 100~MHz, lo que el modo generador cumple por órdenes de
magnitud. La ruta de caracterización no cruza bit a bit: el búfer de
captura se escribe en el dominio del anillo a plena velocidad en una
memoria de 4096 bits y se lee después en el dominio del sistema, con las
dos direcciones del protocolo de acuerdo (handshake) transportadas por
banderas de conmutación. La escritura y la lectura nunca se solapan, así
que no hace falta una FIFO asíncrona, y el tamaño del búfer se deduce de
la Sección~\ref{sec:estimator}: un fragmento de $N + M_{\max}$ bits es
todo lo que necesita el estimador.

\section{Interfaz de registros I2C y secuenciador de órdenes}
\label{sec:i2c}

La interfaz del producto es \ac{I2C}, como en los elementos seguros
comerciales: dos pines, direccionable y el bus que el ESP32 usa para sus
periféricos. El esclavo implementa el protocolo de puntero de registro
con autoincremento, un filtro de pulsos espurios (glitches) de cinco
muestras en ambas líneas, como exige la especificación, y salidas en
drenador abierto. No alarga el reloj (clock stretching): a 400~kHz hay
250 ciclos del sistema por bit, de sobra para cualquier acceso a
registro, y las operaciones largas (carga de semilla, generación) se
lanzan como órdenes y se consultan por sondeo a través de un registro de
estado, que es como se comportan los dispositivos comerciales. El
esclavo se probó con un maestro funcional de bus (bus-functional);
aparecieron dos defectos reales --- el puntero no se autoincrementaba en
las escrituras de varios bytes, y el bit más significativo de cada byte
leído se generaba un flanco de reloj tarde, de modo que todas las
lecturas habrían salido desplazadas --- y se corrigieron.

El mapa de registros (Apéndice~\ref{app:regs}) expone la identificación,
el estado y el control, la configuración del divisor y de la página de
captura, los diagnósticos de los tests de salud, el contador de
resembrados, el reto y su etiqueta, un registro de salida aleatoria
de 256 bits y, \emph{solo mientras está activo un pin de prueba}, la
clave actual y el búfer de captura. El registro de salida aleatoria se
llena mediante su propia orden con una generación separada del DRBG,
nunca con la que produjo la clave: una versión anterior del diseño
dejaba en ese registro los 256 bits de la llamada que generaba la clave,
legibles sin el pin de prueba, y así filtraba la clave. La revisión que
lo detectó se describe en el Capítulo~\ref{ch:ch5_verification}. En la FPGA el pin
de prueba es un interruptor; en un circuito integrado sería un fusible
que se funde en producción, porque exponer la clave o la entropía en
bruto es una vulnerabilidad. Un pequeño secuenciador serializa las
órdenes: la carga de semilla recoge 384 bits en bruto, espera a que los
tests de salud los hayan evaluado todos, aborta si mientras tanto salta
una alarma de salud, e instancia o resiembra el generador; la generación toma 256 bits, guarda los 128 primeros como clave, descarta el resto y recarga las subclaves
de CMAC; una orden aparte llena el registro de salida aleatoria a partir
de una generación independiente; la autenticación ejecuta el CMAC sobre
el bloque de reto.

El motor completo se ejercitó tal como lo hará el ESP32: un maestro
\ac{I2C} funcional de bus lee la identificación, configura el divisor,
carga la semilla, genera, lee la clave en modo de prueba, escribe dos
retos y lee sus etiquetas, vuelve a generar, captura bits en bruto y
verifica que la clave no puede leerse sin el pin de prueba --- dieciséis
comprobaciones. Lo que un banco VHDL no puede saber es si las etiquetas
son \emph{correctas}, porque la clave procede del ruido de los
osciladores en anillo; el banco vuelca la clave, los retos y las
etiquetas a un fichero y un script de Python recalcula los CMAC con el
AES independiente. Coinciden. Es la misma comprobación cruzada que el
ESP32 realiza en hardware con mbedtls.
